\section{Preliminaries: Distribution Matching Distillation}
\label{sec:preliminaries}

Distribution Matching Distillation (DMD) aims to train an efficient student generator $G_\theta$ such that its generated distribution $p_{fake}$ approximates the target real distribution $p_{real}$ defined by a pre-trained diffusion teacher. To ensure consistency between training and inference in few-step synthesis, DMD2~\cite{dmd2} employs \textbf{Backward Simulation} to simulate the inference process. Specifically, starting from pure noise $z \sim \mathcal{N}(0, \mathbf{I})$, the simulation iteratively performs a sequence of denoising and re-noising steps: it predicts a denoised sample $\hat{x}_0$, then adds noise back to reach the next scheduled timestep. This iterative loop continues until a selected timestep $t$ is reached, yielding a simulated clean image that effectively captures the cumulative sampling characteristics of the multi-step inference trajectory. 

This simulated image then serves as the training target, replacing the conventional real data. The generator $G_\theta$ is trained to denoise a noisy version of this simulated sample to get the generated sample $x_g$, with the objective of matching the score functions of the real and fake distributions. The gradient for updating the generator parameters $\theta$ is formulated as:
\begin{equation}
    \label{eq: dmd}
    \nabla_\theta \mathcal{L}_{DMD} = \mathbb{E}_{z, t, \epsilon} \left[ \left( s_{fake}(x_t, t) - s_{real}(x_t, t) \right) \nabla_\theta x_g \right],
\end{equation}
where $x_t$ is the diffused state of the generated sample $x_g$. The target score $s_{real}$ is derived from the frozen teacher, while the fake score $s_{fake}$ is estimated by a critic model updated via standard denoising loss. 
% Additionally, DMD2 \cite{yin2024improved} incorporates a time-conditioned GAN loss to compensate for potential estimation errors in the teacher model. Nevertheless, our empirical observations suggest that this loss leads to significant mode collapse and overfitting in multi-distribution scenarios. Therefore, we depart from the GAN-based approach and exclude it from our design. Beyond the GAN-related issues, directly adapting this framework to non-stationary, customized concepts remains susceptible to vanishing gradients, which we address in the subsequent sections. 
% Additionally, DMD2 \cite{yin2024improved} incorporates a time-conditioned GAN loss to compensate for potential estimation errors in the teacher model. However, adversarial training is inherently unstable, and applying this loss under the limited data constraints of the few-shot personalization setting severely exacerbates mode collapse and overfitting in multi-distribution scenarios. Consequently, this GAN-based objective is excluded from our design. Beyond the challenges associated with adversarial training, directly adapting the distribution matching framework to non-stationary, customized concepts remains susceptible to vanishing gradients, which we address in the subsequent sections.